\documentclass[11pt,a4paper]{ctexart}
\usepackage[a4paper,margin=1.78cm,headheight=15pt]{geometry}
\usepackage{amsmath,amssymb,mathtools,bm}
\usepackage{booktabs,longtable,tabularx,array}
\usepackage{enumitem}
\usepackage[most]{tcolorbox}
\usepackage{tikz}
\usetikzlibrary{arrows.meta,positioning,fit,calc}
\usepackage{xcolor,fancyhdr,hyperref,microtype}
\newcolumntype{Y}{>{\raggedright\arraybackslash}X}
\newcolumntype{L}[1]{>{\raggedright\arraybackslash}p{#1}}
\setlength{\parindent}{2em}\setlength{\parskip}{0.34em}\setlength{\emergencystretch}{3em}
\renewcommand{\arraystretch}{1.22}\setlength{\tabcolsep}{3pt}
\setlist[itemize]{itemsep=0.18em,topsep=0.35em,leftmargin=2em}
\setlist[enumerate]{itemsep=0.18em,topsep=0.35em,leftmargin=2em}
\definecolor{deep}{RGB}{42,58,73}\definecolor{soft}{RGB}{245,247,249}
\definecolor{accent}{RGB}{104,76,55}\definecolor{greenish}{RGB}{57,92,76}
\definecolor{warn}{RGB}{136,68,63}\definecolor{purple}{RGB}{87,72,116}
\hypersetup{colorlinks=true,linkcolor=deep,urlcolor=accent,pdftitle={条件概率、独立性与 Bayes}}
\pagestyle{fancy}\fancyhf{}
\lhead{\small\color{deep}\textbf{数学一 · 概率统计 Topic 02}}
\rhead{\small\color{deep}条件化 · 路径 · 独立 · Bayes}\cfoot{\small\thepage}
\renewcommand{\headrulewidth}{0.4pt}
\newtcolorbox{corebox}[1]{breakable,colback=soft,colframe=deep,title=\textbf{#1},boxrule=0.8pt,arc=2mm}
\newtcolorbox{methodbox}[1]{breakable,colback=white,colframe=greenish,title=\textbf{#1},boxrule=0.7pt,arc=1.5mm}
\newtcolorbox{warnbox}[1]{breakable,colback=white,colframe=warn,title=\textbf{#1},boxrule=0.7pt,arc=1.5mm}
\newtcolorbox{examplebox}[1]{breakable,colback=white,colframe=purple,title=\textbf{#1},boxrule=0.7pt,arc=1.5mm}
\newcommand{\kw}[1]{\textbf{\color{deep}#1}}

\title{\vspace{-1.1cm}\textbf{条件概率、独立性与 Bayes}\\[0.4em]
\large 信息怎样限制世界、合成路径并重分配权重}
\author{数学一 · 概率统计 Topic 02}\date{}

\begin{document}\maketitle

\begin{corebox}{母问题：得到新信息以后，怎样在不破坏概率一致性的前提下重述随机世界？}
已知 $B$ 发生后，$B^c$ 中的结果已被排除；但剩余质量仍须重新归一化。于是
\[
\boxed{\Omega\xrightarrow{\text{observe }B}B\xrightarrow{\text{renormalize}}P(\cdot\mid B)}.
\]
沿多步路径相乘得到联合概率；把隐藏来源的互斥路径相加得到全概率；观察结果后把路径权重重新归一化得到 Bayes。独立性回答反向问题：加入某条信息后，概率是否完全没有变化？
\end{corebox}

\section{Owner 边界与对象}
\begin{itemize}
\item \kw{Owns：}事件层条件概率、乘法公式、完备事件组、全概率、Bayes、事件独立、两两/相互独立、条件独立及事件关联。
\item \kw{Uses：}Topic 01 的 $(\Omega,\mathcal F,P)$、事件代数和互斥可加性。
\item \kw{Outputs：}向 Topic 03/04 提供“信息限制并归一化”的机制；联合密度的支撑计算由 Topic 04 Owns。
\item \kw{Stops：}不把条件概率写成因果干预，也不把 Bayes 的事件后验混成 Topic 08 的参数似然估计。
\end{itemize}

\section{条件概率：在缩小世界中重新称量}

设 $P(B)>0$，定义
\[
\boxed{P(A\mid B)=\frac{P(A\cap B)}{P(B)}}.
\]
分子只保留同时满足 $A,B$ 的结果，分母把新世界 $B$ 的总质量归一化为一。固定 $B$ 后，$A\mapsto P(A\mid B)$ 仍满足非负性、规范性和可列可加性；因此条件概率不是临时技巧，而是一套新的概率测度。

直接得到
\begin{align*}
P(A^c\mid B)&=1-P(A\mid B),\\
A_1\subseteq A_2&\Longrightarrow P(A_1\mid B)\le P(A_2\mid B),\\
P(A_1\cup A_2\mid B)&=P(A_1\mid B)+P(A_2\mid B)-P(A_1\cap A_2\mid B).
\end{align*}
注意补集仍相对于原 $\Omega$，只是条件世界把 $B$ 外部质量视为零。一般 $P(A\mid B)\ne P(B\mid A)$，因为分母分别是 $P(B)$ 与 $P(A)$。

\begin{warnbox}{定义的存在条件}
初等条件概率要求条件事件正概率。连续模型中的非空零概率事件不能直接代入分母；遇到此类问题必须使用更高级的正则条件概率理论，不能假装 $0/0$ 有意义。
\end{warnbox}

\begin{examplebox}{骰子条件化}
骰子结果等可能，$A=\{4,5,6\}$，$B=\{2,4,6\}$。则 $A\cap B=\{4,6\}$，故
\[
P(A\mid B)=\frac{2/6}{3/6}=\frac23.
\]
条件世界只剩 $\{2,4,6\}$；不是把原概率表上的 $P(A)$ 直接改名。
\end{examplebox}

\section{乘法公式：一条路径的质量}

从定义移项得
\[
\boxed{P(A\cap B)=P(A\mid B)P(B)=P(B\mid A)P(A)}.
\]
多事件链式法则为
\[
\boxed{P\!\left(\bigcap_{i=1}^nA_i\right)=P(A_1)\prod_{k=2}^{n}P\!\left(A_k\mid\bigcap_{i=1}^{k-1}A_i\right)},
\]
每个条件都必须有定义。它的结构是“先进入第一阶段，再在已经发生的全部历史下继续前进”，不是把所有概率无条件相乘。

\section{完备事件组与全概率}

若 $B_1,\ldots,B_n$ 两两互斥且 $\bigcup_iB_i=\Omega$，则任何事件 $A$ 都有唯一互斥分解
\[
A=\mathbin{\dot\bigcup}_{i=1}^{n}(A\cap B_i).
\]
因此
\[
\boxed{P(A)=\sum_{i=1}^{n}P(A\cap B_i)=\sum_{i=1}^{n}P(B_i)P(A\mid B_i)}.
\]
这就是全概率公式。它不要求 $A$ 与 $B_i$ 独立；它要求的是来源分类互斥且完备。可数分割时使用相应的可数和；若 $P(B_i)=0$，初等条件概率 $P(A\mid B_i)$ 未必有定义，此时应保留联合概率形式或删去零质量分支。

\subsection{概率树与分割选择}

一棵概率树把每条路径的条件概率相乘，把到达同一终点的互斥路径相加。好的分割满足：每条结果恰好落入一个类别；类别能让条件概率容易计算；类别数量不过多。分割不能因方便而重叠，也不能因粗略而漏掉来源。

\begin{methodbox}{全概率的结果校验}
将 $P(A)$ 看作条件概率的加权平均，必有
\[
\min_iP(A\mid B_i)\le P(A)\le\max_iP(A\mid B_i).
\]
若算出的总概率超出这些界，说明分割、路径或条件方向有误。
\end{methodbox}

\section{Bayes：观察结果后重分配来源权重}

若已经观察到 $A$，希望反推 $B_j$，由条件概率、乘法与全概率得到
\[
\boxed{P(B_j\mid A)=\frac{P(A\mid B_j)P(B_j)}{\displaystyle\sum_{i=1}^{n}P(A\mid B_i)P(B_i)}}.
\]
四个对象必须分开：先验 $P(B_j)$、似然 $P(A\mid B_j)$、证据 $P(A)$、后验 $P(B_j\mid A)$。似然是“假设固定时结果出现的相容性”，不是关于假设自动归一化的概率分布。

比例与赔率形式揭示更新机制：
\[
P(B_j\mid A)\propto P(A\mid B_j)P(B_j),
\]
\[
\boxed{\frac{P(B_j\mid A)}{P(B_k\mid A)}=\frac{P(A\mid B_j)}{P(A\mid B_k)}\frac{P(B_j)}{P(B_k)}}.
\]
后验排序由“似然乘先验”决定。贝叶斯公式只是概率更新，不自动建立因果关系。

\subsection{二元 Bayes 与基础率}

当来源为 $B,B^c$ 时，
\[
P(B\mid A)=\frac{P(A\mid B)P(B)}{P(A\mid B)P(B)+P(A\mid B^c)P(B^c)},
\]
\[
\frac{P(B\mid A)}{P(B^c\mid A)}=\frac{P(A\mid B)}{P(A\mid B^c)}\frac{P(B)}{P(B^c)}.
\]
医学检测中，$P(A\mid B)$ 是灵敏度，$P(A^c\mid B^c)$ 是特异度，$P(B\mid A)$ 是阳性预测值；三者不能互换。

\begin{examplebox}{罕见病的自然频数解释}
患病率 $P(B)=0.001$，灵敏度 $P(A\mid B)=0.99$，特异度 $0.95$，所以假阳性率为 $0.05$。在 $100000$ 人中，真阳性 $99$ 人，假阳性 $4995$ 人，阳性总数 $5094$ 人，故
\[
P(B\mid A)=\frac{99}{5094}\approx1.94\%.
\]
高灵敏度并不意味着阳性后验高；基础率决定健康人群能产生多少假阳性。
\end{examplebox}

\begin{examplebox}{质量追溯}
三工厂供货比例为 $0.30,0.50,0.20$，次品率为 $0.10,0.05,0.15$。次品路径权重分别为 $0.030,0.025,0.030$，故 $P(D)=0.085$，观察次品后来自第三厂的概率为 $0.030/0.085\approx35.3\%$。后验提高来自“先验供货量”和“次品相容性”的共同作用。
\end{examplebox}

\section{独立性：信息没有概率收益}

事件 $A,B$ 独立的定义是
\[
\boxed{P(A\cap B)=P(A)P(B)}.
\]
若 $P(B)>0$，等价于 $P(A\mid B)=P(A)$；若 $0<P(B)<1$，还等价于 $P(A\mid B^c)=P(A)$。独立是概率关系，不是因果关系。

独立事件的补事件仍独立：$A\perp B$ 蕴含 $A^c\perp B$、$A\perp B^c$、$A^c\perp B^c$。概率为 $0$ 或 $1$ 的事件与任意事件独立，但这属于退化独立，不能用条件概率反向解释，因为某些条件概率没有定义。

更一般地，若有限事件组 $A_1,\ldots,A_s,B_1,\ldots,B_t$ \kw{整体相互独立}，那么只用 $A_1,\ldots,A_s$ 做补、有限并交等事件运算得到的事件，与只用 $B_1,\ldots,B_t$ 做相应运算得到的事件仍独立。等价地说，两组所生成的有限事件代数彼此独立。直观上，独立性属于“两组信息互不改变概率”的结构，而不只属于最初写下来的两个事件。这里不能把“所有跨组两两独立”偷换成整体相互独立；后者才足以保证任意组内事件运算后的独立封闭。

正概率独立事件必有正交集概率，因而不能互斥。两个事件既互斥又独立，当且仅当至少一个事件概率为零。

\subsection{两两独立与相互独立}

两两独立只要求任意二元子集满足乘积；相互独立要求任意 $k\ge2$ 个子集都满足
\[
P(A_{i_1}\cap\cdots\cap A_{i_k})=\prod_{r=1}^{k}P(A_{i_r}).
\]
三事件需同时检查三个二交和一个三交。抛两次公平硬币，令 $A=$第一次正面，$B=$第二次正面，$C=$两次相同，则三者两两独立，但
\[
P(A\cap B\cap C)=\frac14\ne\frac18=P(A)P(B)P(C),
\]
因此不相互独立。$n\ge4$ 时，即使所有两两交和全交都符合乘积，也仍可能缺少中间阶数的子族条件。

\begin{examplebox}{真题闸门：1999 Q五 的容斥不能越级}
设 $A,B,C$ 两两独立，$A\cap B\cap C=\varnothing$，且三者概率均为 $p<1/2$。容斥只允许对二重交使用两两独立：
\[
P(A\cup B\cup C)=3p-3p^2+P(ABC)=3p-3p^2.
\]
若该并集为 $9/16$，则 $3p^2-3p+9/16=0$，代数根为 $1/4,3/4$。但概率方程解出的根还必须对应一个真实可实现的原子分解。以“只发生 $A$”这一原子为例，
\[
\begin{aligned}
P(A\cap B^c\cap C^c)
&=P(A)-P(AB)-P(AC)+P(ABC)\\
&=p-2p^2=p(1-2p)\ge0.
\end{aligned}
\]
只要 $p>0$，模型自身就强制 $p\le1/2$，因此 $p=3/4$ 根本不能对应任何满足题设独立结构的概率空间，最终只能取 $p=1/4$。这也说明题面的额外条件 $p<1/2$ 对筛根并非必需；它只是把模型可实现性提前写在了题面上。

\textbf{更一般的校验：}事件概率题把未知量化成代数方程后，候选解必须再通过所有基本原子概率非负、总和为 $1$ 以及题面包含/互斥/独立结构。代数上成立只是必要步骤，概率模型可实现才是最终资格。关键同时有两层：容斥时不能把“两两独立”越级成相互独立；解完方程后也不能把代数根越级成概率答案。
\end{examplebox}

\begin{examplebox}{条件方向比较的共同核心：1998 Q五、2017 Q七}
若 $0<P(A),P(B)<1$，则条件概率方向可以统一还原为交集偏离量：
\[
P(A\mid B)-P(A\mid B^c)
=\frac{P(A\cap B)-P(A)P(B)}{P(B)P(B^c)}.
\]
因此
\[
P(A\mid B)>P(A\mid B^c)
\Longleftrightarrow P(A\cap B)>P(A)P(B)
\Longleftrightarrow P(B\mid A)>P(B\mid A^c).
\]
特别地，若 $P(B\mid A)=P(B\mid A^c)$，全概率先给 $P(B)=P(B\mid A)$，再得 $P(A\cap B)=P(A)P(B)$，所以 $A,B$ 独立。比较条件概率时不要直接交换条件方向；先通分，才能看出它们其实共享同一个联合结构。
\end{examplebox}

\section{关联与条件独立}

定义事件偏离独立的量
\[
\Delta(A,B)=P(A\cap B)-P(A)P(B).
\]
$\Delta>0$ 为正关联，$\Delta<0$ 为负关联，$\Delta=0$ 为独立。若 $0<P(A)<1$，正关联等价于
\[
P(B\mid A)>P(B)>P(B\mid A^c),
\]
负关联则不等号反向。关联只描述概率变化，不是因果结论。

用指示变量 $I_A=\mathbf1_A$，有
\[
E(I_A)=P(A),\quad E(I_AI_B)=P(A\cap B),\quad
\operatorname{Cov}(I_A,I_B)=\Delta(A,B).
\]
因此对事件而言，指示变量协方差为零等价于独立；这条等价不能推广到一般随机变量（一般“不相关”不推出独立）。

若 $P(C)>0$，条件独立定义为
\[
\boxed{P(A\cap B\mid C)=P(A\mid C)P(B\mid C)}.
\]
在相应分母正时，也可写为 $P(A\mid B\cap C)=P(A\mid C)$。无条件独立与给定 $C$ 的条件独立互不蕴含；加入背景信息可能创造依赖，也可能消除依赖。

\begin{methodbox}{条件独立的混合边界：层内乘法不等于边缘乘法}
条件独立只给出每一层的乘法。若 $C$ 有多个可能层 $C_i$，则
\[
P(A\cap B)=\sum_iP(C_i)P(A\mid C_i)P(B\mid C_i),
\]
而 $P(A)P(B)$ 是两个加权和的乘积，两者一般不同。最小反例是 $C\sim\operatorname{Bernoulli}(1/2)$ 且 $A=B=C$：给定 $C$ 后两事件都是常值，故条件独立；但 $P(A\cap B)=1/2\ne1/4=P(A)P(B)$。看到“在每个组内独立”时，必须先分层混合，不能直接删掉边缘协方差。
\end{methodbox}

\begin{examplebox}{真题动作链：2021 Q十六 的状态转移制造相关}
甲、乙两盒各有 2 个红球、2 个白球。先从甲盒取球，记
$X=1$ 表示取到红球；将该球放入乙盒后，再从乙盒取球，记
$Y=1$ 表示取到红球。第一步改变了第二步的组成，因此不能把
$X,Y$ 当作独立指标。先写条件层：
\[
P(X=1)=\frac12,\qquad
P(Y=1\mid X=1)=\frac35,\qquad
P(Y=1\mid X=0)=\frac25.
\]
全概率给出边缘概率
\[
P(Y=1)=\frac12\cdot\frac35+\frac12\cdot\frac25=\frac12.
\]
若题目继续问相关性，必须保留联合路径而不是只看两个边缘：
\[
E(XY)=P(X=1,Y=1)=\frac12\cdot\frac35=\frac3{10},
\]
\[
\operatorname{Cov}(X,Y)=\frac3{10}-\frac12\cdot\frac12=\frac1{20},
\qquad
D(X)=D(Y)=\frac14,
\qquad \rho_{XY}=\frac15.
\]
\kw{模型句：}边缘概率相同只说明混合后“平均位置”相同；只要条件分布随前一阶段状态改变，联合交叉矩就会留下依赖。若目标只问 $P(Y=1)$，在全概率闭合处停止；若问协方差或独立性，继续保留 $P(X,Y)$ 的路径账本。
\end{examplebox}

\begin{examplebox}{2012 真题：指数竞争必须沿“另一者尚未发生”积分}
设 $X\sim\operatorname{Exp}(1)$、$Y\sim\operatorname{Exp}(4)$ 且相互独立，求 $P(X<Y)$。固定 $X=x$ 时，事件 $X<Y$ 要求第二个等待时间仍未结束，因此
\[
P(X<Y)=\int_0^\infty P(Y>x)f_X(x)\,dx
 =\int_0^\infty e^{-4x}e^{-x}\,dx=\frac15.
\]
一般地，若 $X\sim\operatorname{Exp}(\lambda_X)$、$Y\sim\operatorname{Exp}(\lambda_Y)$ 独立，则
\[
P(X<Y)=\int_0^\infty e^{-\lambda_Yx}\lambda_Xe^{-\lambda_Xx}\,dx
 =\frac{\lambda_X}{\lambda_X+\lambda_Y}.
\]
这里的速率决定“先发生”的竞争优势；不能比较均值 $1/\lambda_X$、$1/\lambda_Y$ 后直接写概率，也不能把 $\lambda_X$ 当成概率。\textbf{停止条件：}只问先后概率时，生存函数与密度积分闭合即可；若继续问最小等待时间或胜者条件分布，再转入 P03/P04。
\end{examplebox}

\begin{examplebox}{2022 真题：条件正态先翻译成共享信号模型}
若 $X\sim N(0,1)$，且给定 $X=x$ 时 $Y\mid X=x\sim N(x,1)$，不要把条件均值 $E(Y\mid X)=X$ 直接当成恒等式 $Y=X$。等价的生成表示是
\[
Y=X+\varepsilon,\qquad \varepsilon\sim N(0,1),\quad \varepsilon\perp X.
\]
于是条件层给出的边缘方差为
\[
E(Y)=E[X]=0,\qquad
\operatorname{Var}(Y)=E[\operatorname{Var}(Y\mid X)]+\operatorname{Var}(E[Y\mid X])=1+1=2,
\]
而共享信号留下
\[
\operatorname{Cov}(X,Y)=\operatorname{Cov}(X,X+\varepsilon)=\operatorname{Var}(X)=1,
\qquad
\rho_{XY}=\frac1{\sqrt2}.
\]
所以 $X,Y$ 虽然都是正态边缘，却不独立；独立噪声只说明在给定 $X$ 后的剩余波动，不会消除 $X$ 同时进入 $Y$ 所产生的联合依赖。若题目只问 $Y$ 的边缘矩，在全方差分解闭合处停止；若继续问联合分布、协方差或独立性，必须保留这条共享信号路径并交给 P04/P05 的接口。
\end{examplebox}

\begin{examplebox}{真题分工：2018 Q十四 的独立与互斥各管一段}
题设 $A\perp B$、$A\perp C$、$B\cap C=\varnothing$、$P(A)=P(B)=1/2$，且
\[
P(AC\mid AB\cup C)=1/4.
\]
令 $c=P(C)$。独立性给出 $P(AC)=c/2$、$P(AB)=1/4$；互斥性给出 $AB\cap C=\varnothing$，所以分母为 $1/4+c$。于是
\[
\frac{c/2}{1/4+c}=\frac14\Longrightarrow c=\frac14.
\]
不能把 $A\perp B$ 与 $A\perp C$ 合并成 $B\perp C$；本题正是用互斥而不是独立来处理分母。
\end{examplebox}

\begin{examplebox}{事件指示变量先恢复四格表：2004 Q二十二}
令 $X=\mathbf1_A,Y=\mathbf1_B$。题设 $P(A)=1/4$、$P(B\mid A)=1/3$、$P(A\mid B)=1/2$，先由路径乘法反解
\[
P(AB)=\frac1{12},\qquad P(B)=\frac{P(AB)}{P(A\mid B)}=\frac16.
\]
于是四格联合分布为
\[
P(1,1)=\frac1{12},\quad P(1,0)=\frac16,
\quad P(0,1)=\frac1{12},\quad P(0,0)=\frac23.
\]
再计算
\[
\operatorname{Cov}(X,Y)=\frac1{12}-\frac14\cdot\frac16=\frac1{24},\quad
\rho_{XY}=\frac{1/24}{\sqrt{(3/16)(5/36)}}=\frac1{\sqrt{15}}.
\]
条件概率给的是联合表的约束；只有四格表闭合后，相关性和独立性判断才有依据。
\end{examplebox}

\section{重复试验与可靠性：独立性的结构化应用}

Bernoulli 试验只有成功/失败两种结果，成功概率为 $p$。独立重复 $n$ 次，成功次数 $X\sim\operatorname{Bin}(n,p)$：
\[
P(X=k)=\binom nkp^k(1-p)^{n-k},\qquad
P(X\ge1)=1-(1-p)^n.
\]
“至少 $k$ 次”既可直接求尾和，也可对 $0,\ldots,k-1$ 次使用补事件。

若部件正常事件 $A_i$ 相互独立，串联系统正常需全部正常：
\[
R_{\rm series}=P(\cap_iA_i)=\prod_ip_i.
\]
并联系统只需至少一个正常：
\[
R_{\rm parallel}=1-P(\cap_iA_i^c)=1-\prod_i(1-p_i).
\]
串并联公式的前提不是“看起来像系统”，而是部件事件的独立结构；有共同原因或共享环境时，必须回到条件概率/全概率模型。

\section{统一解题流程与错误边界}

\begin{enumerate}
\item 圈出“已知、在……条件下、观察到、来自哪个来源”等信息词，定义目标事件和条件事件。
\item 检查条件事件或每个分割分支是否具有正概率；零概率分支不要硬套条件概率。
\item 判断方向：原因到结果用路径乘法/全概率，结果到原因用 Bayes。
\item 若用全概率，先验证分类互斥且完备；若用独立，先验证乘积结构，不能凭直觉。
\item 写联合路径 $P(B_i)P(A\mid B_i)$，再求和或归一化；不要直接交换 $P(A\mid B)$ 与 $P(B\mid A)$。
\item 检查结果在 $[0,1]$，后验和为 $1$，加权平均落在条件概率最小值与最大值之间。
\end{enumerate}

\begin{center}
\begin{tabularx}{\linewidth}{L{4.2cm}Y}
\toprule
错误信号 & 纠偏问题 \\
\midrule
$P(A\mid B)$ 与 $P(B\mid A)$ 混用 & 当前分母代表哪个有效世界？ \\
似然直接当后验 & 是否乘了先验并除以全部证据？ \\
分支重叠或遗漏 & 每个结果是否恰好落入一个 $B_i$？ \\
互斥当独立 & 交集为空是集合关系，还是乘积关系？ \\
两两独立当相互独立 & 是否检查了所有非空子族？ \\
关联当因果 & 题目是否真的改变了生成机制？ \\
\bottomrule
\end{tabularx}
\end{center}

\section{可复原的心智模型}
\begin{center}
\begin{tikzpicture}[>=Stealth]
  \tikzset{
    flowstage/.style={draw=deep, fill=white, rounded corners=1.5mm, minimum height=0.72cm, align=center, font=\small\sffamily\bfseries\color{deep}, line width=0.6pt, inner xsep=6pt, inner ysep=3.5pt},
    flowarr/.style={->, >=Stealth, line width=0.75pt, draw=accent}
  }
  \node[flowstage] (n1) at (0,0) {1. 限制世界};
  \node[flowstage, right=0.45cm of n1] (n2) {2. 沿路径相乘};
  \node[flowstage, right=0.45cm of n2] (n3) {3. 跨路径相加};
  \node[flowstage, below=0.6cm of n3] (n4) {4. 按证据归一化};
  \node[flowstage, left=0.45cm of n4] (n5) {5. 检查信息是否改变概率};

  \draw[flowarr] (n1) -- (n2);
  \draw[flowarr] (n2) -- (n3);
  \draw[flowarr] (n3) -- (n4);
  \draw[flowarr] (n4) -- (n5);
\end{tikzpicture}
\end{center}
条件概率是“限制并归一化”；全概率是“把隐藏路径合成”；Bayes 是“观察终点后重新分配路径”；独立性是“信息没有改变权重”。

\begin{examplebox}{条件概率为 1 先翻译成包含关系：2006 真题}
若 $P(B)>0$ 且 $P(A\mid B)=1$，则
\[
P(B\setminus A)=P(B)-P(AB)=P(B)\bigl[1-P(A\mid B)\bigr]=0.
\]
因此 $B\subseteq A$ 几乎处处，$A\cup B$ 与 $A$ 只差一个零概率集合，故 $P(A\cup B)=P(A)$。这不是“条件概率等于 1 就能把两个事件写成相等”，而是概率意义下的包含。
\end{examplebox}

\begin{examplebox}{早期真题：来源分割与 Bayes 分母（1987）}
随机等概率选三箱，白球率分别为 $1/5,1/2,5/8$。先算证据概率
\[
P(W)=\frac13\left(\frac15+\frac12+\frac58\right)=\frac{53}{120},
\]
再得
\[
P(\text{箱2}\mid W)=\frac{(1/3)(1/2)}{53/120}=\frac{20}{53}.
\]
先验权重、似然和证据分母必须分开；似然不是后验。
\end{examplebox}

\end{document}
